\oldpage[IV-4]{86}

  \item[\textup{c$'$)}] $h$ est lisse au point $y$ et il existe un voisinage ouvert $U$ de $y$ dans $Y$ tel que $j|U:U\to X$ soit une immersion quasi-régulière \textup{(autrement dit (16.9.8), il existe un voisinage $U$ de $y$ dans $Y$ dans lequel $\sh{N}_{Y/X}$ est un $\sh{O}_Y$-Module localement libre et l'homomorphisme canonique $\mathbf{S}^{\bullet}_{\sh{O}_Y}(\sh{N}_{Y/X})\to\mathscr{G}r^{\bullet}(j)$ est bijectif)}.
\end{enumerate}
\end{theorem}

Pour démontrer l'équivalence de \emph{a)} et \emph{b)}, on peut se borner au cas où $S=\Spec(A)$ et $X=\Spec(B)$ sont affines, avec $Y=\Spec(B/\mathfrak J)$, où $\mathfrak J$ est un idéal de type fini de $B$, et $B$ est une $A$-algèbre lisse (17.3.2, (ii)). Il suffit alors d'appliquer le critère jacobien (\textbf{0}, 22.6.1) ainsi que (\textbf{0}, 19.1.14), compte tenu de ce que $\varOmega^1_{X/S}$ est un $\sh{O}_X$-Module localement libre dans un voisinage de $x$ et $\sh{N}_{Y/X}$ un $\sh{O}_Y$-Module de type fini (16.1.6).

En second lieu, prouvons que \emph{a)} entraîne \emph{c')}. On peut encore se borner au cas où $S$, $X$, $Y$ sont affines, $B$ et $B/\mathfrak J$ des $A$-algèbres lisses, et il résulte alors de (17.7.9) et (8.10.5, (iv)) qu'il existe un sous-anneau noethérien $A'$ de $A$, une $A'$-algèbre de type fini $B'$ et un idéal $\mathfrak J'$ de $B'$ tels que $B=B'\otimes_{A'}A$, $\mathfrak J=\mathfrak J'B$ et que $B'$ et $B'/\mathfrak J'$ soient des $A'$-algèbres lisses. Notons que, d'après (\textbf{0}, 19.3.8), $B'$ est encore une $A'$-algèbre formellement lisse lorsqu'on munit $A'$ de la topologie discrète et $B'$ de la topologie $\mathfrak J'$-préadique. Par suite (\textbf{0}, 19.5.4), $\mathfrak J'/\mathfrak J'^2$ est un $(B'/\mathfrak J')$-module projectif et l'homomorphisme canonique $\mathbf{S}^{\bullet}_{B'/\mathfrak J'}(\mathfrak J'/\mathfrak J'^2)\to\operatorname{gr}_{\mathfrak J'}^{\bullet}(B')$ est bijectif. En d'autres termes (16.9.8), l'immersion $\Spec(B'/\mathfrak J')\to\Spec(B')$ est quasi-régulière, donc régulière puisque $B'$ est noethérien (16.9.10). Mais comme par hypothèse $B'/\mathfrak J'$ est un $A'$-module plat (17.5.1) et $B'$ une $A'$-algèbre de présentation finie, on peut appliquer (11.3.8) en remplaçant $\sh{F}$ par $\widetilde{\mathfrak J}'$, et on voit donc (par changement de base) que $j$ est une immersion régulière, et \emph{a fortiori} quasi-régulière. Le fait que $B$ est une $A$-algèbre de présentation finie et un $A$-module plat montre d'ailleurs, en vertu de (11.3.8), que les conditions \emph{c')} et \emph{c)} sont équivalentes, et sont stables par changement de base.

Montrons enfin que \emph{c)} entraîne \emph{a)}. Du fait que, dans (11.3.8), la condition \emph{b)} entraîne \emph{c)}, on voit déjà que si \emph{c)} est vérifiée, $f$ est \emph{plat} au point $x$; en vertu de (17.5.1), tout revient donc à voir que sous l'hypothèse \emph{c)}, $f^{-1}(s)$ est lisse sur $\boldsymbol{k}(s)$ au point $x$, en posant $s=h(x)$; comme on a remarqué que la condition \emph{c)} est stable par changement de base, on voit qu'on est ramené au cas où $S=\Spec(k)$ est le spectre d'un corps, et compte tenu de (17.7.1, (ii)) on peut supposer que $k$ est \emph{algébriquement clos}. Il suffit alors de prouver que $\sh{O}_{X,x}$ est un anneau \emph{régulier} (17.5.1). Or, par hypothèse on a $\sh{O}_{Y,y}=\sh{O}_{X,x}/\sh{J}_x$, où $\sh{J}_x$ est un idéal engendré par une suite $\sh{O}_{X,x}$-régulière $(t_i)$, et $\sh{O}_{Y,y}$ est un anneau \emph{régulier}; comme les $t_i$ font partie d'un système de paramètres pour $\sh{O}_{X,x}$ (\textbf{0}, 16.4.1), la conclusion résulte de (\textbf{0}, 17.1.7).

\begin{corollary}[17.12.2]
\label{IV.17.12.2}
Avec les notations de \textup{(17.12.1)}, supposons que $f$ soit lisse au point $x$, et soit $Y$ un sous-préschéma fermé de $X$ défini par un Idéal $\sh{J}$ de $\sh{O}_X$, et $S$-lisse au point $x$. Soient $(g_i)_{1\leq i\leq r}$ des sections de $\sh{J}$ au-dessus de $X$. Les conditions suivantes sont équivalentes:
\begin{enumerate}
  \item[\textup{a)}] Les $(g_i)_x$ forment un système de générateurs du $\sh{O}_{X,x}$-module $\sh{J}_x$ dont le nombre d'éléments est le plus petit possible.
  \item[\textup{b)}] Si $g_i'$ est l'image canonique de $g_i$ dans $\Gamma(X,\sh{J}/\sh{J}^2)$, les $(g_i')_x$ forment une base du $(\sh{O}_{X,x}/\sh{J}_x)$-module $\sh{J}_x/\sh{J}_x^2$.

